\section{结论}

本实验利用电容-电压法（C--V法）对五个${\rm p}^+{\rm n}$单边突变结二极管进行了系统测量，结合锁定放大器微弱信号检测技术，成功实现了轻掺杂一侧杂质浓度分布的测定。

实验首先验证了锁定放大器从噪声中提取微弱信号的能力。在不同时间常量$T_c$（1--1000 ms）下的噪声测试表明，输出电压的波动$\sigma(R_{1\omega})$与$1/\sqrt{T_c}$成正比，符合低通滤波器等效噪声带宽$\Delta f_{\rm N}=1/(4RC)$的理论预期，验证了锁定放大器通过压缩带宽实现信噪比改善的基本原理。

固定电容校准实验确认了锁定放大器输出电压$V_{\rm out}$与待测电容$C_{\rm x}$之间良好的线性关系，为C--V曲线电容轴的定量标定提供了可靠依据。

通过对五个二极管的C--V特性测量和相位分析，筛选出D$_2$为适用于C--V法分析的合格单边突变结二极管。D$_2$的$1/C^{2}$--$V_{\rm R}$曲线线性度优异，相位随偏压变化幅度小且接近参考相位，表明其反向漏电流小、串联电阻低。而D$_3$、D$_4$、D$_5$由于存在较大的反向漏电流，相位显著偏离参考值，不适用于定量杂质分布分析。这一筛选过程表明，相位$(\phi_{\rm v}-\phi_0)$与偏压$V_{\rm R}$的关系是判断pn结质量的有效判据。

利用D$_2$的$C$--$V_{\rm R}$数据计算得到了轻掺杂一侧的杂质浓度分布$N(w)$。在空间电荷区深度$w$约0.2--0.6 $\mu$m范围内，$N(w)$分布在$10^{15}$--$10^{16}\ {\rm cm}^{-3}$量级，符合扩散工艺制备的浅pn结典型掺杂水平。通过$1/C^{2}$--$V_{\rm R}$直线段的外推，获得了接触电势差$V_{\rm D}$和均匀掺杂段的掺杂浓度$N_{\rm D}$。

本实验结果表明，C--V法是一种非破坏性、快速有效的半导体杂质分布测量方法。正确选择合格的二极管样品（以相位接近参考值为判据）是保证测量结果可靠性的关键前提。锁定放大器作为微弱信号检测的核心仪器，其时间常量的合理选取需在信噪比与测量速度之间折中，实验中$T_c = 100$--300 ms可获得较好的综合性能。
